\chapter{The Head of Desk's Job}\label{ch:fm:the-head-of-desks-job}

The budget meeting in December sets next year's revenue target at this year's result plus ten per cent, and leaves the desk's risk limits where
they were. This year's result was exactly what the desk expects to make; its revenue has a Sharpe ratio of 1.3. With those limits the desk meets
the new target in 45 years out of a hundred. The head of desk has been handed a plan that fails more often than it succeeds, before a single
trade; to meet it three years in four, the desk would have to run more than twice its risk. This chapter is about the numbers a head of desk
agrees to in December and runs by for the rest of the year.

\section{The budget: revenue, costs and the plan}

\begin{definition}[Revenue budget, allocated cost]\label{def:fm:the-head-of-desks-job:budget}
A desk's \emph{revenue budget}\index{revenue budget} is the revenue it commits to produce in the coming year, agreed with the firm's management
and used to plan its costs, its pay pool and its capital. An \emph{allocated cost}\index{allocated cost} is a share of the firm's shared costs,
technology, premises, control functions, market data, charged to the desk by an allocation key (headcount, revenue, usage).
\end{definition}

A desk's plan has three parts: revenue by business, costs (its own direct costs and those allocated to it), and the capital and limits it will use
(chapter 7). The chapter's desk (illustrative, \$ millions a year) runs three businesses: client flow, expected revenue 30 and volatility 25;
inventory trading, 15 and 20; structured solutions, 10 and 12; correlated 0.3, 0.2 and 0.4. Its expected revenue is 55, its volatility 42.2
and its Sharpe ratio 1.30. Direct costs of 18 and \omterm{def:fm:the-head-of-desks-job:budget}{allocated costs} of 10 leave, after variable pay of 20\% of a positive result, an expected
profit of 20.5, and a 26\% chance that revenue does not cover costs.

\begin{proposition}[What a target above expectation costs]\label{prop:fm:the-head-of-desks-job:target}
Let annual revenue be normal with mean $\mu$ and volatility $\sigma$, Sharpe ratio $S=\mu/\sigma$, and let the budget be $k\mu$. The probability of
meeting it is
\[
P=\Phi\bigl(S\,(1-k)\bigr).
\]
For $k>1$ it is below one half and falls as the Sharpe ratio rises. Scaling every position by $\lambda$ (mean and volatility both scale) meets the
budget with probability $q$ if and only if $\lambda\,(\mu-z_q\sigma)=k\mu$, $z_q=\Phi^{-1}(q)$, which requires $S>z_q$.
\end{proposition}

\begin{proof}
$P(\text{revenue}\ge k\mu)=1-\Phi((k\mu-\mu)/\sigma)=\Phi(S(1-k))$, decreasing in $S$ when $1-k<0$. After scaling, revenue is normal with mean
$\lambda\mu$ and volatility $\lambda\sigma$: $P=q$ iff $(\lambda\mu-k\mu)/(\lambda\sigma)=z_q$, i.e.\ $\lambda(\mu-z_q\sigma)=k\mu$, which has a
positive solution only if $\mu>z_q\sigma$.
\end{proof}

The counter-intuitive half of the proposition is the useful one: a steadier desk is \emph{less} likely to beat a stretch target, because its
results cluster closer to their mean (\cref{fig:fm:the-head-of-desks-job:pmeet}). At $k=1.1$ and $S=1.3$ the desk meets its budget with
probability 44.8\% (45.0\% in 20\,000 simulated years); to meet it with probability 75\% it must scale every position by 2.28. A budget that sits
above the desk's expected revenue does not ask for effort: it asks for risk.

\begin{omfigure}
\begin{tikzpicture}
\begin{axis}[width=10.5cm,height=5.4cm,xlabel={\small budget as a multiple of expected revenue $k$},ylabel={\small probability of meeting it (\%)},
  xmin=0.5,xmax=1.5,ymin=0,ymax=100,tick label style={font=\footnotesize},grid=major,grid style={gray!20},table/col sep=comma,
  legend cell align=left,legend style={font=\scriptsize,at={(0.5,-0.3)},anchor=north,/tikz/every even column/.append style={column sep=8pt}},legend columns=4]
\addplot[omDef,thick] table[x=k,y=s05]{figdata/desk/06-the-head-of-desks-job/pmeet.csv};
\addplot[omProp,thick,dashed] table[x=k,y=s10]{figdata/desk/06-the-head-of-desks-job/pmeet.csv};
\addplot[omThm,thick,dotted] table[x=k,y=s15]{figdata/desk/06-the-head-of-desks-job/pmeet.csv};
\addplot[gray,thick,dashdotted] table[x=k,y=s20]{figdata/desk/06-the-head-of-desks-job/pmeet.csv};
\draw[black,thin] (axis cs:1,0) -- (axis cs:1,100);
\legend{Sharpe 0.5,Sharpe 1.0,Sharpe 1.5,Sharpe 2.0}
\end{axis}
\end{tikzpicture}

\omcaption{The probability that a desk's annual revenue meets a budget set at $k$ times its expected revenue, $\Phi(S(1-k))$, for four Sharpe
ratios. Every curve passes through one half at $k=1$; above it, the steadier the desk, the less likely it beats the budget. Data:
\texttt{fm\_head.p\_meet\_closed}.}\label{fig:fm:the-head-of-desks-job:pmeet}
\end{omfigure}

\begin{method}[Negotiating a budget]\label{met:fm:the-head-of-desks-job:budget}
\begin{enumerate}
\item Estimate each business's expected revenue and volatility from its own history, on current limits, not from last year's result alone.
\item Compute the desk's expected revenue, volatility and Sharpe ratio, and the probability of meeting the proposed budget
(\cref{prop:fm:the-head-of-desks-job:target}).
\item If the budget is above expectation, state the risk scale it requires and ask for the limits that go with it, or for a budget at
expectation.
\item Agree the costs allocated to the desk and the key that allocates them; check that the desk covers its full costs in its expected year.
\end{enumerate}
\end{method}

\section{Risk appetite: from a loss the firm can bear to limits}

\begin{definition}[Risk appetite, risk appetite statement]\label{def:fm:the-head-of-desks-job:appetite}
A firm's \emph{risk appetite}\index{risk appetite} is the aggregate level and types of risk it is willing to take, within the maximum it could
take given its capital, liquidity and obligations, to achieve its strategy and business plan. Its \emph{risk appetite
statement}\index{risk appetite statement} is the written form: qualitative statements and quantitative measures relative to earnings,
capital, risk and liquidity, which its risk limits then allocate to business lines, desks and risk types.
\end{definition}

The definitions follow the Financial Stability Board's principles of 2013, which also name the upper bound, \emph{risk capacity}, the most
risk a firm can take before breaching its regulatory capital and liquidity constraints. A statement such as ``the firm accepts a loss of at most
\$20 million in a one-in-twenty year from this desk'' becomes limits in three steps (\cref{fig:fm:the-head-of-desks-job:cascade}).

\begin{omfigure}
\begin{tikzpicture}[font=\footnotesize,>=stealth,box/.style={draw,rounded corners,minimum height=0.7cm,align=center,fill=omDef!10,
  minimum width=3.2cm}]
\node[box,fill=omProp!12] (lt) at (0,2.6) {loss tolerance\\{\scriptsize \$20m in a 1-in-20 year}};
\node[box] (vb) at (0,1.2) {volatility budget\\{\scriptsize $(\mu+L)/z_{0.95}=45.6$}};
\node[box] (var) at (-3.6,-0.4) {daily VaR limit\\{\scriptsize $z_{0.99}\,\sigma/\sqrt{252}=6.7$}};
\node[box] (st1) at (0.2,-0.4) {stop: client flow\\{\scriptsize 8.8}};
\node[box] (st2) at (3.8,-0.4) {stops: inventory 7.0,\\{\scriptsize structured 4.2}};
\draw[->] (lt) -- (vb); \draw[->] (vb) -- (var); \draw[->] (vb) -- (st1); \draw[->] (vb) -- (st2);
\node[font=\scriptsize,align=left,anchor=west] at (2.1,1.2) {desk's volatility today 42.2:\\it may scale up by 1.08};
\end{tikzpicture}

\omcaption{The risk-appetite cascade for the chapter's desk: a loss tolerance becomes a volatility budget, which becomes a daily value-at-risk limit
and a stop for each business, shared by volatility. \$ millions; illustrative desk. Model: \texttt{firm.deskplan.cascade}.}\label{fig:fm:the-head-of-desks-job:cascade}
\end{omfigure}

With expected revenue $\mu=55$, a loss of at most $L=20$ at 95\% confidence allows a volatility of $(55+20)/1.645=45.6$: the desk, at 42.2, may
scale up by 8\%. At 99\% one day, the corresponding value at risk (One Quant Book 6, chapter 21) is $2.326\times45.6/\sqrt{252}=6.7$. The \$20
million of tolerance, shared by each business's volatility, gives stops of 8.8, 7.0 and 4.2. The cascade is the risk hierarchy of One Quant Book
6, chapter 29, read from the top: the limits are the \omterm{def:fm:the-head-of-desks-job:appetite}{risk appetite statement} in numbers.

\omcode{code/firm/deskplan/firm_deskplan.py}{73}{103}{The risk scale a budget requires, and the cascade from a loss tolerance to a volatility
budget, a daily VaR limit and stops.}\label{lst:fm:the-head-of-desks-job:cascade}

\begin{remark}[The budget and the appetite must agree]\label{rem:fm:the-head-of-desks-job:agree}
The desk's \omterm{def:fm:the-head-of-desks-job:appetite}{risk appetite} allows a volatility of 45.6; its budget, to be met three years in four, needs 2.28 times 42.2, or 96. A plan that sets
one without the other asks the head of desk to break one of them. The planning round is where the firm chooses: a lower budget, a larger appetite,
or a better Sharpe ratio, and only the first two can be chosen in December.
\end{remark}

\section{Product scope and new-product approval}

\begin{definition}[Product scope, new-product approval]\label{def:fm:the-head-of-desks-job:scope}
A desk's \emph{product scope}\index{product scope} is the list of instruments, markets and client types it is authorised to trade. \emph{New-product
approval}\index{new-product approval} is the process by which a firm reviews a new or modified product, market or activity before a desk may
trade it: its risks, pricing and valuation, booking and settlement, legal and tax treatment, capital and compliance, each signed off by the
function responsible.
\end{definition}

A new product is where a desk's risks are least understood and its controls least tested: the pricing model is new, the booking may be manual, the
settlement process may not exist. US bank supervisors' guidance expects banks to manage the risks of ``new, modified, or expanded products and
services'' with appropriate approvals before a new activity starts and with controls to identify, measure, monitor and report its risks. For the
head of desk, the approval process is the price of growth: every function that signs is one that will support the product when it goes wrong.
\omterm{def:fm:the-head-of-desks-job:scope}{Product scope} also fixes what the desk may \emph{not} do; a trader who finds a profitable trade outside it has found a request for approval, not
a trade.

\begin{method}[Taking a new product through approval]\label{met:fm:the-head-of-desks-job:npa}
\begin{enumerate}
\item Write the product's description: payoff, clients, expected volume, the desk's risk and its hedges.
\item For each function (risk, product control, operations, legal, tax, compliance, technology) list what it must build or check, and get its
sign-off with conditions.
\item Start with limits on size and a review date; move the product to business as usual only after the first review.
\end{enumerate}
\end{method}

\section{The rhythm: the day, the week, the year}

A head of desk runs to a calendar that repeats. The contents vary by firm; the structure is common, because it follows the controls:

\begin{center}
\small
\begin{tabular}{@{}lp{10cm}@{}}
\toprule
when & what the head of desk does\\
\midrule
every morning & read the previous day's P\&L and its explanation (One Quant Book 1, chapter 7), limit usage and breaks; the morning meeting on
markets, positions and client flow\\
during the day & approve trades above delegated limits; watch intraday risk and stops; client calls\\
every evening & sign off the day's P\&L with product control (chapter 15); check positions against limits\\
every week & review each business against plan, its risk and its attribution; hiring and projects\\
every month & plan against actual (\cref{fig:fm:the-head-of-desks-job:fan}); the risk committee's review (chapter 12)\\
every year & budget and limits; pay pool and its allocation (chapter 10); \omterm{def:fm:the-head-of-desks-job:scope}{product scope}\\
\bottomrule
\end{tabular}
\end{center}

The monthly review compares cumulative revenue with the plan. A plan spread evenly over the year is a line; the desk's revenue is a random path
around its expectation, and the band it can wander in without meaning anything widens with the square root of time
(\cref{fig:fm:the-head-of-desks-job:fan}). In month three, a desk exactly on its expected path is already \$1.4 million behind a budget ten per
cent above it, well inside its noise; a head of desk who reacts to that gap by adding risk is reacting to noise.

\begin{omfigure}
\begin{tikzpicture}
\begin{axis}[width=10.5cm,height=5.6cm,xlabel={\small month},ylabel={\small cumulative revenue (\$ million)},xmin=0,xmax=12,ymin=-60,ymax=140,
  xtick={0,2,4,6,8,10,12},tick label style={font=\footnotesize},grid=major,grid style={gray!20},table/col sep=comma,
  legend cell align=left,legend style={font=\scriptsize,at={(0.5,-0.3)},anchor=north,/tikz/every even column/.append style={column sep=8pt}},legend columns=4]
\addplot[name path=a,draw=none,forget plot] table[x=month,y=p5]{figdata/desk/06-the-head-of-desks-job/fan.csv};
\addplot[name path=b,draw=none,forget plot] table[x=month,y=p95]{figdata/desk/06-the-head-of-desks-job/fan.csv};
\addplot[omDef!15] fill between[of=a and b];
\addplot[name path=c,draw=none,forget plot] table[x=month,y=p25]{figdata/desk/06-the-head-of-desks-job/fan.csv};
\addplot[name path=d,draw=none,forget plot] table[x=month,y=p75]{figdata/desk/06-the-head-of-desks-job/fan.csv};
\addplot[omDef!40] fill between[of=c and d];
\addplot[omDef,thick] table[x=month,y=p50]{figdata/desk/06-the-head-of-desks-job/fan.csv};
\addplot[omThm,thick,dashed] table[x=month,y=budget]{figdata/desk/06-the-head-of-desks-job/fan.csv};
\legend{5th--95th percentile,25th--75th percentile,median,budget}
\end{axis}
\end{tikzpicture}

\omcaption{The chapter's desk: percentiles of cumulative revenue by month over 5\,000 simulated years, against a budget of 1.1 times expected
revenue spread evenly. By December the budget sits just above the median path. Data: \texttt{fm\_head.fan}.}\label{fig:fm:the-head-of-desks-job:fan}
\end{omfigure}

\section{Tutorial: next year's plan}\label{tut:fm:the-head-of-desks-job:tut}

\begin{tutorial}
\textbf{Goal.} Build a desk's plan, measure how likely it is to be met, and derive its limits from a loss tolerance. \textbf{End state:}
\cref{fig:fm:the-head-of-desks-job:fan} and the numbers of \cref{rem:fm:the-head-of-desks-job:agree}.
\begin{tutsteps}
\item \textbf{The plan.} \texttt{fm\_head.PLAN} is a \texttt{firm.deskplan.Plan} of three businesses, their correlations, costs and the
budget; \texttt{plan\_stats} gives expected revenue 55, volatility 42.2, Sharpe ratio 1.30.
\item \textbf{The odds.} \texttt{p\_meet} gives 44.8\%; \texttt{simulate} confirms it on 20\,000 years and returns the maximum drawdowns (median
\$30.0 million, 90th percentile \$50.4 million).
\item \textbf{The risk the budget needs.} \texttt{risk\_multiple(PLAN, 0.75)} gives 2.28 (\cref{lst:fm:the-head-of-desks-job:cascade}).
\item \textbf{The cascade.} \texttt{cascade(PLAN, 20)} gives the volatility budget, the daily VaR limit and the stops.
\item \textbf{The year.} \texttt{fm\_head.fan()} gives the monthly percentiles of \cref{fig:fm:the-head-of-desks-job:fan}.
\end{tutsteps}
\end{tutorial}

\textbf{What to change next.} Raise the correlation between flow and inventory to 0.8 and redo the cascade; set the budget at expectation and
compare the drawdowns the desk must live with.

\section{Build: the desk plan}\label{bld:fm:the-head-of-desks-job:deskplan}

\begin{build}
\bfield{Purpose} The head of desk's numbers in one place: the plan, its odds, the risk it needs and the limits the firm's appetite allows.
\bfield{Interface} \texttt{firm.deskplan}: \texttt{Business(name, mu, sigma)}, \texttt{Plan(businesses, corr, direct\_cost, allocated\_cost,
var\_pay, budget)}; \texttt{plan\_stats}, \texttt{p\_meet}, \texttt{risk\_multiple(plan, prob)}, \texttt{simulate(plan, n, rng)},
\texttt{cascade(plan, loss\_tolerance, conf, var\_conf)}, \texttt{plan\_vs\_actual(plan, actual\_monthly)}.
\bfield{Rules} Revenue is normal over the year (the chapter says so where it matters); a budget that needs a Sharpe ratio above $z_q$ returns an
infinite risk multiple; stops add up to the loss tolerance.
\bfield{Acceptance tests} \texttt{code/firm/deskplan/tests/}: the closed form against simulation; the risk multiple solves its equation and
is infinite when the Sharpe ratio is too low; the cascade's volatility budget and stops; the plan-against-actual rows.
\bfield{Stretch} Fat-tailed revenue (a Student-$t$) and its effect on the cascade; stops by business from each business's contribution to the
desk's volatility (chapter 7) instead of its stand-alone volatility.
\end{build}

\begin{omsources}
\item Financial Stability Board, \emph{Principles for an Effective Risk Appetite Framework}, 18 November 2013.
\item Office of the Comptroller of the Currency, Bulletin 2017-43, \emph{New, Modified, or Expanded Bank Products and Services: Risk Management
Principles}, 20 October 2017.
\end{omsources}

\section{Exercises}

\begin{exercise}[$\star$]\label{exo:fm:the-head-of-desks-job:1}
A desk with a Sharpe ratio of 1 is given a budget at 1.1 times its expected revenue. What is its probability of meeting it?
\end{exercise}

\begin{exercise}[$\star$]\label{exo:fm:the-head-of-desks-job:2}
The chapter's desk expects 55 with volatility 42.2 and has costs of 28. What is the probability that revenue does not cover its costs?
\end{exercise}

\begin{exercise}[$\star$]\label{exo:fm:the-head-of-desks-job:3}
Compute the volatility budget of a desk expecting 55 whose firm accepts a loss of at most 20 in a one-in-twenty year, and its daily 99\% VaR limit.
\end{exercise}

\begin{exercise}[$\star\star$]\label{exo:fm:the-head-of-desks-job:4}
Show that a desk of Sharpe ratio 1.3 cannot meet any budget with probability 95\% by scaling its positions, whatever the budget.
\end{exercise}

\begin{exercise}[$\star\star$]\label{exo:fm:the-head-of-desks-job:5}
From \cref{fig:fm:the-head-of-desks-job:pmeet}, at what budget multiple does a desk of Sharpe ratio 2 have the same probability of success as a desk
of Sharpe ratio 0.5 at $k=1.5$?
\end{exercise}

\begin{exercise}[$\star\star$]\label{exo:fm:the-head-of-desks-job:6}
In month three the desk of \cref{fig:fm:the-head-of-desks-job:fan} is exactly on its expected path. How far behind the budget is it, and how many
standard deviations of its three-month revenue is that?
\end{exercise}

\begin{exercise}[$\star\star\star$]\label{exo:fm:the-head-of-desks-job:7}
\emph{Coding.} Rerun \texttt{firm.deskplan.risk\_multiple} for the chapter's desk with the budget at 1.2 times expectation and a target
probability of 75\%. What scale is needed, and what daily VaR would it imply?
\end{exercise}

\begin{exercise}[$\star\star\star$]\label{exo:fm:the-head-of-desks-job:8}
\emph{Find the flaw.} ``We missed budget two years running although our Sharpe ratio is 1.3. The desk is underperforming.''
\end{exercise}

\section{Problem: Plus Ten Per Cent}

\begin{problem}[{Weekend problem --- plus ten per cent}]\label{pb:fm:the-head-of-desks-job:1}
A head of desk is handed next year's budget, ten per cent above this year's result, with unchanged limits. She has a week to answer.

\medskip\noindent\textbf{Part I --- The plan.}
\begin{enumerate}
\item Define a \omterm{def:fm:the-head-of-desks-job:budget}{revenue budget} and an \omterm{def:fm:the-head-of-desks-job:budget}{allocated cost}.
\item Give the desk's expected revenue, volatility and Sharpe ratio from its three businesses.
\item Give its expected profit after costs and variable pay, and the probability that revenue does not cover costs.
\item What is the budget, in \$ million?
\end{enumerate}

\medskip\noindent\textbf{Part II --- The odds.}
\begin{enumerate}[resume]
\item State and prove \cref{prop:fm:the-head-of-desks-job:target}.
\item Give the probability of meeting the budget, in closed form and simulated.
\item Why is a steadier desk less likely to beat a stretch budget?
\item Give the median and 90th percentile of the desk's maximum drawdown in a year.
\item What risk scale meets the budget three years in four?
\end{enumerate}

\medskip\noindent\textbf{Part III --- The appetite.}
\begin{enumerate}[resume]
\item Define \omterm{def:fm:the-head-of-desks-job:appetite}{risk appetite} and a \omterm{def:fm:the-head-of-desks-job:appetite}{risk appetite statement}, and name the upper bound on \omterm{def:fm:the-head-of-desks-job:appetite}{risk appetite}.
\item Turn a \$20 million loss tolerance at 95\% into a volatility budget.
\item Give the daily 99\% VaR limit and the stops by business.
\item By how much may the desk scale up under its appetite?
\item Compare the risk the appetite allows with the risk the budget needs.
\end{enumerate}

\medskip\noindent\textbf{Part IV --- The answer.}
\begin{enumerate}[resume]
\item Define \omterm{def:fm:the-head-of-desks-job:scope}{product scope} and \omterm{def:fm:the-head-of-desks-job:scope}{new-product approval}, and say why a new product is where controls are weakest.
\item In month three the desk is on its expected path. How far behind the budget is it, in \$ million and in standard deviations?
\item What should the head of desk not do in response?
\item Which of the plan's inputs can be changed in December, and which cannot?
\item State the \emph{named result}: the probability of meeting a budget of $k$ times expected revenue at Sharpe ratio $S$, its value for this
desk, and the risk scale needed to meet it three years in four.
\item In two sentences, write her answer to the budget.
\end{enumerate}
\end{problem}

\section{Interview questions}

\begin{interviewq}[$\star$ \iqroles{trader}]\label{iq:fm:the-head-of-desks-job:1}
Your desk's Sharpe ratio is 1. What is the probability it makes less than its expected revenue in a year? Less than zero?
\end{interviewq}

\begin{interviewq}[$\star$ \iqroles{risk}]\label{iq:fm:the-head-of-desks-job:2}
What is a \omterm{def:fm:the-head-of-desks-job:appetite}{risk appetite statement}, and how does it become a desk's limits?
\end{interviewq}

\begin{interviewq}[$\star\star$ \iqroles{trader, researcher}]\label{iq:fm:the-head-of-desks-job:3}
Management raises your budget 20\% and keeps your limits. What does that do to your odds, and what would you ask for?
\end{interviewq}

\begin{interviewq}[$\star\star$ \iqroles{trader}]\label{iq:fm:the-head-of-desks-job:4}
You are three months into the year and behind a budget spread evenly. When should that worry you?
\end{interviewq}

\begin{interviewq}[$\star\star$ \iqroles{bank, risk}]\label{iq:fm:the-head-of-desks-job:5}
A trader wants to trade a product the desk has never traded. What has to happen first, and why?
\end{interviewq}

\begin{interviewq}[$\star\star\star$ \iqroles{researcher, risk}]\label{iq:fm:the-head-of-desks-job:6}
Why can a higher Sharpe ratio lower the probability of meeting a budget? When does it raise it?
\end{interviewq}
